\section{Attribute Grammar: Synthesized dan Inherited}
\label{sec:ch14-attribute}

\subsection{Atribut}

\begin{definisi}[Attribute grammar]
Attribute grammar adalah tata bahasa bebas konteks yang setiap
simbolnya memiliki atribut dan setiap produksinya memiliki aturan
semantik untuk menghitung atribut tersebut.
\end{definisi}

Atribut dapat berupa tipe, nilai, alamat, kode, atau informasi lain.
Setiap atribut memiliki domain nilai dan aturan evaluasi.

\subsection{Atribut Synthesized}

Atribut synthesized dihitung dari atribut anak. Contoh klasik adalah
atribut \code{type} pada simpul ekspresi:
\[
E.type = \text{rule}(E_1.type, E_2.type).
\]
Nilai mengalir dari daun ke akar. SDD yang hanya memakai atribut
synthesized disebut \eng{S-attributed}.

\subsection{Atribut Inherited}

Atribut inherited dihitung dari orang tua atau saudara. Contohnya
adalah atribut \code{env} yang membawa lingkungan simbol ke anak:
\[
\text{stmt}.env = \text{block}.env.
\]
Nilai mengalir dari akar ke daun atau antar-saudara. SDD yang
memakai inherited disebut \eng{L-attributed} jika ketergantungannya
mengikuti urutan kiri-ke-kanan.

\begin{figure}[H]
\centering
\begin{tikzpicture}[node distance=8mm and 14mm]
  \node[phase, text width=3.0cm] (root) {\(E\)\\type = float};
  \node[phase, text width=2.4cm, below left=of root] (left)
        {\(E_1\)\\type = int};
  \node[phase, text width=2.4cm, below right=of root] (right)
        {\(E_2\)\\type = float};
  \draw[arrow] (left) -- node[left, font=\small] {synthesized} (root);
  \draw[arrow] (right) -- node[right, font=\small] {synthesized} (root);
\end{tikzpicture}
\caption{Atribut synthesized mengalir dari anak ke orang tua}
\label{fig:ch14-synthesized}
\end{figure}

\subsection{Dependency Graph}

Ketergantungan antar-atribut membentuk graf berarah. Evaluasi atribut
harus mengikuti urutan topologis graf tersebut. Jika graf memiliki
siklus, evaluasi tidak dapat dilakukan tanpa iterasi atau
pendekatan khusus.

\begin{warningbox}
Atribut inherited dapat menimbulkan siklus jika aturan saling
bergantung. Periksa graf ketergantungan sebelum mengimplementasikan
evaluasi.
\end{warningbox}
